%%
%% Test: Index System — v1.1
%% Stage: index refactor (stages 1-4)
%% 
%% Purpose: Verify index system with xindy + persian-variant2:
%%   - \index entries are sorted correctly (پ, چ, ژ, گ, ک)
%%   - \idxbold, \idxitalic work
%%   - \idxsee, \idxseealso work
%%   - \idxsub, \idxsubsub work
%%   - \printindex works
%%   - Index appears in ToC (intoc)
%% 

\documentclass{matinbook}

\title{تست سیستم نمایه — نسخه ۱.۱}
\author{تیم MatinBook}
\date{۱۴۰۵}

\booktitle{تست نمایه}

\begin{document}

\maketitle

\tableofcontents

\chapter{بررسی سیستم نمایه}
\label{chap:test-index}

این سند، سیستم نمایه \texttt{mb-index} را بررسی می‌کند.
هدف اصلی، تست مرتب‌سازی صحیح حروف فارسی (پ، چ، ژ، گ، ک) است.

%----------------------------------------
\section{حروف الفبای فارسی}
\label{sec:persian-letters}

در این بخش، کلماتی با حروف خاص فارسی را نمایه می‌کنیم:

\textbf{حرف پ:}
پایتون\index{پایتون}، برنامه\index{برنامه}، پردازش\index{پردازش}

\textbf{حرف چ:}
چاپ\index{چاپ}، چندجمله‌ای\index{چندجمله‌ای}، چرخه\index{چرخه}

\textbf{حرف ژ:}
ژنتیک\index{ژنتیک}، ژرف\index{ژرف}

\textbf{حرف گ:}
گسسته\index{گسسته}، گراف\index{گراف}، گویا\index{گویا}

\textbf{حرف ک:}
کامپیوتر\index{کامپیوتر}، کتابخانه\index{کتابخانه}، کلاس\index{کلاس}

\textbf{حروف عادی:}
الگوریتم\index{الگوریتم}، برنامه‌نویسی\index{برنامه‌نویسی}، ریاضیات\index{ریاضیات}

%----------------------------------------
\section{نمایه‌های تودرتو}
\label{sec:nested-index}

برنامه‌نویسی\index{برنامه‌نویسی!پایتون}\index{برنامه‌نویسی!جاوا}\index{برنامه‌نویسی!سی‌پلاس‌پلاس}

الگوریتم\index{الگوریتم!مرتب‌سازی}\index{الگوریتم!جستجو}

ریاضیات\index{ریاضیات!جبر}\index{ریاضیات!آنالیز}\index{ریاضیات!هندسه}

%----------------------------------------
\section{نمایه‌های bold و italic}
\label{sec:bold-italic}

مفهوم مهم\idxbold{مفهوم مهم}

اصطلاح لاتین\idxitalic{الگوریتم مرتب‌سازی}

%----------------------------------------
\section{ارجاعات see و seealso}
\label{sec:see-also}

\idxsee{پایتون}{زبان برنامه‌نویسی}

\idxseealso{برنامه‌نویسی}{الگوریتم}

%----------------------------------------
\section{نمایه‌های تودرتوی عمیق}
\label{sec:deep-nested}

\idxsubsub{ریاضیات}{جبر}{گروه}

\idxsubsub{ریاضیات}{جبر}{حلقه}

\idxsubsub{علوم کامپیوتر}{الگوریتم}{مرتب‌سازی}

%----------------------------------------
\section{تکرار برای range}
\label{sec:repeated}

الگوریتم\index{الگوریتم} یکی از مفاهیم اساسی است.

در اینجا دوباره به الگوریتم\index{الگوریتم} اشاره می‌کنیم.

و باز هم الگوریتم\index{الگوریتم} را تکرار می‌کنیم.

%----------------------------------------
\section{نتیجه‌گیری}
\label{sec:conclusion}

اگر این سند بدون خطا کامپایل شود:

\begin{itemize}
    \item ✅ \texttt{xindy} با \texttt{persian-variant2} کار می‌کند
    \item ✅ حروف پ، چ، ژ، گ، ک درست مرتب می‌شوند
    \item ✅ \texttt{\textbackslash idxbold} و \texttt{\textbackslash idxitalic} کار می‌کنند
    \item ✅ \texttt{\textbackslash idxsee} و \texttt{\textbackslash idxseealso} کار می‌کنند
    \item ✅ \texttt{\textbackslash idxsub} و \texttt{\textbackslash idxsubsub} کار می‌کنند
    \item ✅ \texttt{\textbackslash printindex} کار می‌کند
    \item ✅ نمایه در فهرست مطالب نمایش داده می‌شود
\end{itemize}

\textbf{وضعیت تست:} قبول

%----------------------------------------
% Print the index
%----------------------------------------
\printindex

\end{document}
